\section{The MFS for Maxwell's equations}
\label{sec:method}

\subsection{The scattering problem}

We consider the computation of the electromagnetic field in the unbounded region exterior to a perfect electric conductor (PEC). Here, the conductor occupies a bounded domain $\Omega\subset\mathbb{R}^3$ with surface $\Gamma$, and we work in the frequency domain with time dependence $\rme^{-\rmi\omega t}$ and wavenumber $k$. The incident field is the plane wave
\begin{equation}
  \Einc(\bx) = \bp\,\rme^{\rmi k\,\bd\cdot\bx},
  \label{eq:planewave}
\end{equation}
where $\bp$ is the polarisation vector and $\bd$ the propagation direction. Note that $\bp$ and $\bd$ are orthogonal unit vectors, which is essential for electromagnetic waves in free space. The scattered field $\bE$ satisfies \cite{kirsch2015,colton2013}
\begin{equation}
  \nabla\times\nabla\times\bE - k^2\,\bE = 0 \quad\text{in } \mathbb{R}^3\setminus\overline{\Omega},
  \label{eq:maxwell}
\end{equation}
together with the Silver--M\"uller radiation condition $\lim_{r\to\infty} r\bigl((\nabla\times\bE)\times\hat\bx - \rmi k\,\bE\bigr) = 0$, where $r = \|\bx\|$ and $\hat\bx = \bx/r$, which selects outgoing waves at infinity. On a perfect conductor, the physical requirement is that the tangential component of the total electric field vanishes on $\Gamma$,
\begin{equation}
  \bt\cdot\bigl(\bE + \Einc\bigr) = 0 \quad\text{on } \Gamma
  \label{eq:pec}
\end{equation}
for every tangent vector $\bt$. Equivalently, the currents induced on the surface produce a scattered field that exactly cancels the tangential component of the incident field. Far from the body, the scattered field behaves as an outgoing spherical wave, $\bE(\bx) = \rme^{\rmi kr}\,\bF(\hat\bx)/r + O(r^{-2})$, where $\bF(\hat\bx)$ is the far-field pattern, and the radar cross-section in the observation direction $\hat\bx$ is $4\pi|\bF(\hat\bx)|^2$.

\subsection{The fundamental solution}

Let $\Phi(r) = \rme^{\rmi kr}/(4\pi r)$ with $r = \|\bx-\by\|$ be the outgoing fundamental solution of the scalar Helmholtz equation $(\Delta + k^2)\Phi = -\delta(\bx-\by)$. To lift this scalar solution to a divergence-free vector solution of \eqref{eq:maxwell} away from $\by$, we follow \cite{osipov2017} and apply the operator $I + k^{-2}\nabla\nabla$, since $\nabla\times(\nabla\nabla\Phi\,\balpha) = 0$ and $\Delta\Phi = -k^2\Phi$ for $\bx\ne\by$. Here, $I$ is the unit dyadic, which acts as the $3\times3$ identity on vectors, and $\nabla\nabla\Phi$ is the Hessian of $\Phi$, the $3\times3$ matrix of its second partial derivatives. Writing $\rhat = (\bx-\by)/r$ and noting that the outer product $\rhat\rhat^{\top}$ projects onto the radial direction, we obtain
\begin{equation}
\begin{split}
  \bG(\bx,\by) &= \Bigl(I + \frac{1}{k^2}\nabla\nabla\Bigr)\Phi(r)\\
  &= \Phi(r)\Biggl(\Bigl(1+\frac{\rmi}{kr}-\frac{1}{(kr)^2}\Bigr)I-\Bigl(1+\frac{3\rmi}{kr}-\frac{3}{(kr)^2}\Bigr)\rhat\rhat^{\top}\Biggr).
\end{split}
  \label{eq:dyadic}
\end{equation}
Because this Green's function is a $3\times3$ matrix that connects vector sources to vector fields, it is often referred to as the \emph{dyadic Green's function}. For a vector $\balpha\in\mathbb{C}^3$, the field $\bG(\bx,\by)\,\balpha$ is that of an electric dipole at $\by$ with moment $\balpha$; this field satisfies both \eqref{eq:maxwell} away from $\by$ and the radiation condition. Taking $(1/k)\nabla_{\bx}\times(\bG(\bx,\by)\,\bbeta)$ eliminates the Hessian term and gives a second family of electric-field solutions, radiated by magnetic dipoles at $\by$ with moment $\bbeta\in\mathbb{C}^3$,
\begin{equation}
  \frac{1}{k}\,\nabla_{\bx}\Phi(r)\times\bbeta
  = \Bigl(\rmi-\frac{1}{kr}\Bigr)\Phi(r)\;\rhat\times\bbeta .
  \label{eq:magnetic}
\end{equation}
For the smooth bodies of \autoref{sec:smooth} the electric dipoles are sufficient, whereas bodies with edges need both families, for reasons given in \autoref{sec:skeleton}.

\subsection{Formulation}

We choose $N$ source points $\by_1,\dots,\by_N$ inside the conductor and approximate the scattered field by
\begin{equation}
  \bE_N(\bx) = \sum_{j=1}^{N}\Bigl[\bG(\bx,\by_j)\,\balpha_j + \frac{1}{k}\nabla_{\bx}\Phi(\|\bx-\by_j\|)\times\bbeta_j\Bigr],
  \label{eq:mfs}
\end{equation}
where each coefficient $\balpha_j$ or $\bbeta_j$ is a 3-component vector. As usual, the formulation can be simplified in the implementation by absorbing the constant factor $1/(4\pi)$ of $\Phi$ into the coefficients $\balpha_j$ and $\bbeta_j$, and the coefficient norms reported below follow this convention. Since the approximation satisfies \eqref{eq:maxwell} and the radiation condition for any choice of the coefficients, only the boundary condition \eqref{eq:pec} has to be enforced. We do this at $M$ boundary points $\bx_1,\dots,\bx_M$ on $\Gamma$. At each boundary point an orthonormal tangent basis $\bt_i^{1},\bt_i^{2}$ is computed, and the coefficients are chosen to minimise
\begin{equation}
  \sum_{i=1}^{M} w_i \sum_{s=1}^{2}\bigl|\bt_i^{s}\cdot\bigl(\bE_N(\bx_i)+\Einc(\bx_i)\bigr)\bigr|^2,
  \label{eq:leastsquares}
\end{equation}
where the weight $w_i$ is the surface area that the point represents (modified near the edges of the cube, \autoref{sec:cubepoints}). Since each coefficient is a vector and each dipole field is given by a $3\times3$ matrix, the resulting system is not a simple scalar one: it is composed of $3\times3$ blocks, one for each pair of a boundary point and a dipole, and after projection onto the tangent plane each block is reduced to a $2\times3$ block. Scaling the two rows at $\bx_i$ by $\sqrt{w_i}$ turns \eqref{eq:leastsquares} into a linear least-squares problem $\|A\balpha - \mathbf{b}\|_2$ with $2M$ equations and $3N$ unknowns for electric dipoles (or $6N$ when magnetic dipoles are included). We always take more equations than unknowns, so that the boundary condition is also controlled between the boundary points.

A well-documented disadvantage of the MFS is that this system of equations is often ill-conditioned \cite{fairweather1998,barnett2008}. For the sphere and the torus we therefore solve it with orthogonal factorisations \cite{golub2013}: QR for the sphere, and the singular value decomposition for the small systems of the torus. For the layouts of the cube and the almond, in contrast, the matrices are well conditioned once their columns are scaled to unit norm, with condition numbers below $10^{6}$ in almost all cases (the exception is the cube at three wavelengths for $J\ge9$, where the QR fallback below is used). Because squaring a condition number below $10^{6}$ keeps $\operatorname{cond}(A^{*}A)$ safely within double precision, in these cases we solve the normal equations $A^{*}A\mathbf{c} = A^{*}\mathbf{b}$ via a Cholesky factorisation $A^{*}A = R^{*}R$ and apply two steps of iterative refinement, $\mathbf{c}\leftarrow\mathbf{c} + (R^{*}R)^{-1}A^{*}(\mathbf{b}-A\mathbf{c})$, following the corrected semi-normal equations of Bj\"orck \cite{bjorck1987}. This costs about half as much as QR, and the fields of the cube agree with those of a QR solve to $10^{-13}$. If the Cholesky factorisation finds the matrix too ill-conditioned, the solver falls back to QR. In all cases the cost is $O(MN^2)$ operations, and the systems have at most a few thousand unknowns once the symmetry of the body has been used.

\subsection{Measuring the error}
\label{sec:errors}

For the sphere, the Mie series \cite{mie1908,bohren1998} gives the exact scattered field, and we report the relative error
\begin{equation}
  \varepsilon = \Bigl(\sum\nolimits_{l}|\bE_N(\mathbf{z}_l)-\bE(\mathbf{z}_l)|^2\Bigr)^{1/2}\Big/\Bigl(\sum\nolimits_{l}|\bE(\mathbf{z}_l)|^2\Bigr)^{1/2}
  \label{eq:fielderror}
\end{equation}
at points $\mathbf{z}_l$ on a sphere that encloses the body: $400$ points on the sphere of radius $2$ for the unit sphere and the unit cube ($1500$ points at radius $1.5$ in \autoref{fig:sphere}), $200$ points on the sphere of radius $2.5$ for the torus, and $400$ points on the sphere of radius $0.3$\,m for the almond. For the non-spherical shapes no closed-form solution exists, so we evaluate \eqref{eq:fielderror} with a reference field in place of $\bE$, always stating which. We also report the boundary residual, computed as the square root of \eqref{eq:leastsquares} relative to its value for $\Einc$ alone on the cube and the almond, and as the corresponding root-mean-square tangential residual at $1500$ test points on $\Gamma$ for the sphere and the torus.

We also use two further checks. The first is the optical theorem. Since a perfect conductor absorbs no energy, the total power removed from the incident wave (determined from the forward amplitude $\bF(\bd)$, with $\bar\bp$ the complex conjugate of $\bp$) equals the total power scattered over the unit sphere $S^2$, so the extinction and scattering cross-sections
\begin{equation}
  \sext = \frac{4\pi}{k}\,\operatorname{Im}\bigl(\bar\bp\cdot\bF(\bd)\bigr), \qquad
  \ssca = \int_{S^2}|\bF(\hat\bx)|^2\,\mathrm{d}s(\hat\bx)
  \label{eq:optical}
\end{equation}
are equal for the exact solution. We evaluate $\ssca$ with a Gauss--Legendre rule in the cosine of the polar angle and the trapezoidal rule in the azimuth ($40\times80$ points, or $48\times96$ on the torus and the almond). The quantity $|\sext-\ssca|/\ssca$, referred to below as the optical mismatch, is then a reference-free error indicator. Note, however, that it is not a bound, as an example in \autoref{sec:cubepoints} shows.

The second check is an independent solver. We use the boundary element library Bempp-cl \cite{betcke2021,scroggs2017}, which discretises the surface current in the electric field integral equation with lowest-order divergence-conforming Rao--Wilton--Glisson (RWG) elements \cite{rao1982} on a triangular mesh. It serves both as a reference and as a yardstick for the computing time, and its settings are given in \ref{app:bempp}. All timings in this paper were measured on one machine, with two cores of an Intel Xeon Platinum 8481C processor at $2.2$\,GHz and $8$\,GB of memory, and they include assembly, solution and the evaluation of the field.
